% 核心结论：升高线电压与提高功率因数都通过减小线电流降低输电损耗
% 图原型：纯几何／对比。物理网络按本题节点、端口和符号构建。
\documentclass[tikz,border=7pt]{standalone}
\usepackage{fontspec,xeCJK,amsmath,amssymb}
\IfFontExistsTF{Times New Roman}{\setmainfont{Times New Roman}}{\setmainfont{TeX Gyre Termes}}
\IfFontExistsTF{Heiti SC}{\setCJKmainfont[AutoFakeBold=1.5]{Heiti SC}}{\setCJKmainfont[AutoFakeBold=1.5]{FandolHei-Regular}}
\usetikzlibrary{arrows.meta,calc,angles,quotes}
\definecolor{ink}{HTML}{243745}
\definecolor{blue}{HTML}{275DAB}
\definecolor{green}{HTML}{177C63}
\definecolor{red}{HTML}{BD4935}
\definecolor{purple}{HTML}{715393}
\begin{document}
\begin{tikzpicture}[font=\fontsize{9.5}{12}\selectfont,text=ink,
 v/.style={line width=1.1pt,-{Stealth[length=2.4mm,width=1.5mm]}},
 axis/.style={ink!65,line width=.45pt,-{Stealth[length=1.7mm,width=1.1mm]}},
 aux/.style={ink!40,densely dashed,line width=.5pt},
 lab/.style={fill=white,inner sep=2pt},
 card/.style={draw=blue!25,fill=blue!3,rounded corners=3pt,text width=6cm,align=left,inner sep=10pt}]
\path[use as bounding box] (-8,-3.35) rectangle (8,3.65);
\node[font=\fontsize{10.5}{13}\selectfont\bfseries,text width=15cm,align=center] at (0.00000,3.22000) {M6-H48　升高线电压与提高功率因数都通过减小线电流降低输电损耗};
\begin{scope}[shift={(-4,0)}]
\draw[ink,thick] (-2.80000,1.20000) -- (-1.80000,1.20000);
\draw[ink,thick] (-1.80000,1.20000) -- (-0.82000,1.20000);
\draw[ink,thick] (0.02000,1.20000) -- (1.00000,1.20000);
\begin{scope}[shift={(-0.40000,1.20000)},rotate=0.0]
\draw[ink,thick,fill=white] (-0.42,-.14) rectangle (0.42,.14);
\end{scope}
\node[lab] at (-0.40000,1.62000) {$r$};
\draw[ink,thick] (1.00000,1.20000) -- (2.50000,1.20000);
\node[] at (-3.00000,1.20000) {$A$};
\draw[ink,thick] (-2.80000,0.00000) -- (-1.80000,0.00000);
\draw[ink,thick] (-1.80000,0.00000) -- (-0.82000,0.00000);
\draw[ink,thick] (0.02000,0.00000) -- (1.00000,0.00000);
\begin{scope}[shift={(-0.40000,0.00000)},rotate=0.0]
\draw[ink,thick,fill=white] (-0.42,-.14) rectangle (0.42,.14);
\end{scope}
\node[lab] at (-0.40000,0.42000) {$r$};
\draw[ink,thick] (1.00000,0.00000) -- (2.50000,0.00000);
\node[] at (-3.00000,0.00000) {$B$};
\draw[ink,thick] (-2.80000,-1.20000) -- (-1.80000,-1.20000);
\draw[ink,thick] (-1.80000,-1.20000) -- (-0.82000,-1.20000);
\draw[ink,thick] (0.02000,-1.20000) -- (1.00000,-1.20000);
\begin{scope}[shift={(-0.40000,-1.20000)},rotate=0.0]
\draw[ink,thick,fill=white] (-0.42,-.14) rectangle (0.42,.14);
\end{scope}
\node[lab] at (-0.40000,-0.78000) {$r$};
\draw[ink,thick] (1.00000,-1.20000) -- (2.50000,-1.20000);
\node[] at (-3.00000,-1.20000) {$C$};
\node[] at (0.00000,2.20000) {三相电源 → 线路 → 平衡负载};
\node[] at (0.00000,-2.20000) {$P=\sqrt3U_LI_L\cos\phi$};
\end{scope}
\begin{scope}[shift={(4,0)}]
\node[card] at (0.00000,1.15000) {$I_L=P/(\sqrt3U_L\cos\phi)$};
\end{scope}
\begin{scope}[shift={(4,0)}]
\node[card] at (0.00000,-1.15000) {$P_{loss}=3I_L^2r$};
\end{scope}
\node[text width=15cm,align=center] at (0.00000,-3.00000) {对称正相序与平衡负载；空间磁场矢量和复平面相量分别解释。};

\end{tikzpicture}\end{document}